% (a) Dirichlet函数 1_Q 在[0,1]：有理点取1、无理点取0，两条点线示意稠密；小区间[0.4,0.6]上 sup=1、inf=0。
% (b) Cantor集示性函数 1_C：用第3步的8个保留区间（每段长1/27）示意 C，C外取0；不连续点恰为 C，测度为0。
\begin{tikzpicture}[x=40mm,y=16mm,font=\fontsize{8}{10}\selectfont,
  ax/.style={draw=FigureStroke,line width=.75pt,-{Stealth[length=4pt]}},
  note/.style={font=\FigureNoteFont\fontsize{7.5}{9}\selectfont,text=FigureMuted}]
\begin{scope}
  \fill[FigureYellowFill] (0.4,-0.08) rectangle (0.6,1.12);
  \draw[ax] (-0.03,0)--(1.1,0) node[right,text=FigureInk]{$x$};
  \draw[ax] (0,-0.05)--(0,1.35);
  \draw[FigureStroke,line width=.6pt] (-0.015,1)--(0,1) node[left=1pt,text=FigureInk]{$1$};
  \node[left=1pt,text=FigureInk] at (0,0) {$0$};
  \draw[FigureStroke,line width=.6pt] (1,0)--(1,-0.05) node[below,text=FigureInk]{$1$};
  \foreach \i in {0,...,100}{\fill[FigureBlue] (\i/100,1) circle (0.45pt);}
  \foreach \i in {0,...,100}{\fill[FigureRed] (\i/100+0.005,0.025) circle (0.45pt);}
  \node[text=FigureBlue,anchor=south] at (0.2,1.03) {有理点};
  \node[text=FigureRed,anchor=south] at (0.82,0.05) {无理点};
  \node[note,anchor=north] at (0.5,-0.12) {小区间上 $\sup=1$，$\inf=0$};
  \node[text=FigureInk] at (0.5,1.55) {(a) $\mathbf 1_{\mathbb Q}$};
\end{scope}
\begin{scope}[xshift=52mm]
  \draw[ax] (-0.03,0)--(1.1,0) node[right,text=FigureInk]{$x$};
  \draw[ax] (0,-0.05)--(0,1.35);
  \draw[FigureStroke,line width=.6pt] (-0.015,1)--(0,1) node[left=1pt,text=FigureInk]{$1$};
  \node[left=1pt,text=FigureInk] at (0,0) {$0$};
  \draw[FigureStroke,line width=.6pt] (1,0)--(1,-0.05) node[below,text=FigureInk]{$1$};
  \foreach \a/\b in {1/2,3/6,7/8,9/18,19/20,21/24,25/26}{\draw[FigureRed,line width=1.6pt] (\a/27,0.025)--(\b/27,0.025);}
  \foreach \a in {0,2,6,8,18,20,24,26}{\draw[FigureBlue,line width=1.6pt] (\a/27,1)--(\a/27+1/27,1);}
  \node[note,anchor=north] at (0.5,-0.12) {不连续点恰为 $C$，$\lambda(C)=0$};
  \node[text=FigureInk] at (0.5,1.55) {(b) $\mathbf 1_C$};
\end{scope}
\end{tikzpicture}
